\documentclass[11pt]{article}

\usepackage[margin=1in]{geometry}
\usepackage{amsmath,amssymb,amsfonts,amsthm}
\usepackage{booktabs}
\usepackage{hyperref}
\usepackage{listings}
\usepackage{xcolor}
\usepackage{enumitem}

\hypersetup{colorlinks=true,linkcolor=blue,urlcolor=blue,citecolor=blue}

\lstset{
  basicstyle=\ttfamily\small,
  breaklines=true,
  columns=fullflexible,
  frame=single,
  showstringspaces=false
}

\newtheorem{proposition}{Proposition}

\title{A Paper-Style Technical Note for \texttt{codex\_heat1d\_qutip.py}\\
CV-DV LCHS Simulation of the 1D Heat Equation}
\author{}
\date{\today}

\begin{document}
\maketitle

\begin{abstract}
This note documents, in paper format, the exact model and numerical pipeline
implemented in \texttt{codex\_heat1d\_qutip.py}. The code simulates a
continuous-variable/discrete-variable (CV-DV) linear-combination-of-Hamiltonian-
simulation (LCHS) construction for a 1D heat-equation toy problem with four
interior spatial points (two DV qubits). We provide the mathematical setup,
map each equation to code components, define all reported metrics, and give
command-level reproduction instructions.
The implementation is a dense-matrix SciPy/QuTiP correctness study, not a
hardware resource estimate.
\end{abstract}

\section{Problem Statement}
We consider the heat equation
\begin{equation}
\frac{\partial u(x,t)}{\partial t}=\alpha\frac{\partial^2u(x,t)}{\partial x^2},
\end{equation}
with homogeneous Dirichlet boundary conditions and initial profile $u(x,0)=f(x)$.
Using a second-order finite difference stencil with four interior points,
we obtain
\begin{equation}
\frac{d}{dt}\mathbf{u}(t)=-A\mathbf{u}(t),
\qquad
A=\frac{\alpha}{h^2}
\begin{bmatrix}
2 & -1 & 0 & 0 \\
-1 & 2 & -1 & 0 \\
0 & -1 & 2 & -1 \\
0 & 0 & -1 & 2
\end{bmatrix}.
\label{eq:heatA}
\end{equation}
The classical reference used throughout is
\begin{equation}
\mathbf{u}_{\mathrm{ref}}(t)=e^{-At}\mathbf{u}_0.
\label{eq:classical}
\end{equation}

\paragraph{Pauli verification.}
The script also verifies the two-qubit decomposition
\begin{equation}
A=\frac{\alpha}{h^2}\left(2II-IX-\frac12XX-\frac12YY\right),
\end{equation}
and reports $\lVert A_{\mathrm{direct}}-A_{\mathrm{pauli}}\rVert_F$.

\section{CV-DV LCHS Map Used in the Code}
For this symmetric diffusion generator, the Cartesian split gives $L=A$, $H=0$.
The simulated hybrid map is
\begin{equation}
K(t)=\langle\phi|\exp\bigl(-it\,\hat{x}\otimes A\bigr)|\psi\rangle
\in \mathbb{C}^{4\times4},
\label{eq:Kmap}
\end{equation}
where $\hat{x}$ acts on the oscillator space.

The target LCHS kernel branch is
\begin{equation}
g_\beta(k)=\frac{\exp\!\left(-(1+ik)^\beta\right)}{C_\beta(1-ik)},
\qquad C_\beta=2\pi e^{-2^\beta},
\qquad \beta\in(0,1).
\label{eq:gkernel}
\end{equation}

The oscillator states are parameterized as
\begin{equation}
|\phi\rangle=S(r)|0\rangle,
\qquad
|\psi\rangle=S(r')\sum_{n=0}^{N_c-1}C_n|n\rangle,
\end{equation}
with coefficients from a Gauss--Hermite projection,
\begin{equation}
C_n\propto\int
H_n\!\left(\frac{k}{e^{r'}}\right)
\exp(-\gamma k^2)
\,g_\beta(k)\,dk,
\qquad
\gamma=e^{-2r'}-e^{-2r}>0.
\label{eq:Cn}
\end{equation}
The implementation normalizes $\{C_n\}$ in $\ell_2$.

\section{Exact Mapping to \texttt{codex\_heat1d\_qutip.py}}
\subsection{Core computational blocks}
\begin{itemize}[itemsep=2pt]
  \item \texttt{heat\_matrix()}, \texttt{heat\_matrix\_from\_paulis()}: Eq.~\eqref{eq:heatA} and Pauli check.
  \item \texttt{kernel\_g\_beta()}: Eq.~\eqref{eq:gkernel}.
  \item \texttt{lchs\_coefficients()}: Eq.~\eqref{eq:Cn} with Gauss--Hermite quadrature.
  \item \texttt{prepare\_lchs\_states()}: builds $|\psi\rangle,|\phi\rangle$ in QuTiP.
  \item \texttt{effective\_lchs\_map()}: evaluates Eq.~\eqref{eq:Kmap}.
  \item \texttt{classical\_map()}: computes Eq.~\eqref{eq:classical} via \texttt{scipy.linalg.expm}.
\end{itemize}

\subsection{Position operator convention}
The script supports
\begin{equation}
\hat{x}=\frac{a+a^\dagger}{2}
\quad\text{or}\quad
\hat{x}=\frac{a+a^\dagger}{\sqrt{2}},
\end{equation}
with default \texttt{half} and an empirical multiplier
\texttt{position\_scale=0.8}.

\subsection{Fock-space parameters}
Two cutoffs are exposed:
\begin{itemize}[itemsep=2pt]
  \item \texttt{n-fock}: oscillator Hilbert-space truncation for simulation operators,
  \item \texttt{n-coeff}: number of projected coefficients $C_n$ explicitly computed before zero-padding.
\end{itemize}
By default \texttt{n-coeff=0} in CLI is interpreted as
\texttt{n-coeff = n-fock}. This was added to improve numerical stability when
high-order Hermite terms become ill-conditioned at large $n$.

\section{Metrics and Their Meaning}
Given observed map $K$ and reference map $R=e^{-At}$, the script fits a global
complex scale
\begin{equation}
s_* = \frac{\langle R,K\rangle_F}{\langle R,R\rangle_F}
\end{equation}
and reports
\begin{equation}
\varepsilon_{\mathrm{map}}=\frac{\|K-s_*R\|_F}{\|s_*R\|_F}.
\end{equation}

For a selected input $\mathbf{u}_0$:
\begin{equation}
\mathbf{u}_{\mathrm{cv}}=K\mathbf{u}_0,
\qquad
\mathbf{u}_{\mathrm{ref}}=R\mathbf{u}_0.
\end{equation}
It reports:
\begin{enumerate}[itemsep=2pt]
  \item best-scale vector error (same formula in vector form),
  \item shape fidelity
  $F=|\langle\hat{u}_{\mathrm{cv}},\hat{u}_{\mathrm{ref}}\rangle|^2$,
  \item eta-calibrated error using $\eta=\langle\phi|\psi\rangle$.
\end{enumerate}

\begin{proposition}[Why fidelity and eta-error can disagree]
High fidelity with high eta-calibrated error is possible because fidelity is
scale-insensitive (normalized state direction), while eta-calibrated error is
amplitude-sensitive with a fixed normalization model.
\end{proposition}

\section{Default Numerical Setting and Baseline Output}
Default configuration:
\begin{center}
\begin{tabular}{ll}
\toprule
Parameter & Default \\
\midrule
$\alpha$ & 1.0 \\
$h$ & 1.0 \\
$t$ & 1.0 \\
\texttt{n-fock} & 64 \\
\texttt{n-coeff} & 64 (via \texttt{n-coeff=0} rule) \\
$r$ & 8.0 \\
$r'$ & 4.0 \\
$\beta$ & 0.9 \\
\texttt{n-quad} & 240 \\
initial DV state & \texttt{basis01} \\
$\hat{x}$ convention & \texttt{half} \\
\texttt{position-scale} & 0.8 \\
\bottomrule
\end{tabular}
\end{center}

Default run output (current script):
\begin{itemize}[itemsep=2pt]
  \item map best-scale error: \texttt{0.036052}
  \item vector best-scale error: \texttt{0.023673}
  \item vector shape fidelity: \texttt{0.999440}
  \item eta-calibrated vector error: \texttt{0.361785}
\end{itemize}

\section{Trajectory and Visualization Protocol}
The script supports:
\begin{itemize}[itemsep=2pt]
  \item trajectory CSV/JSON export,
  \item CHO-style component plots for one to four DV components,
  \item separate normalized-error analysis scripts (provided below).
\end{itemize}

\paragraph{Generated files in this workspace.}
Typical outputs are:
\begin{itemize}[itemsep=2pt]
  \item \texttt{results\_bosonic/codex\_heat1d\_qutip\_metrics.json}
  \item \texttt{results\_bosonic/codex\_heat1d\_qutip\_trajectory.csv}
  \item \texttt{results\_bosonic/codex\_heat1d\_qutip\_CHOstyle\_all4.png}
  \item \texttt{results\_bosonic/codex\_heat1d\_qutip\_normalized\_error\_timeseries.csv}
\end{itemize}

\section{Reproducibility Commands}
Run from repository root.

\subsection{Baseline run}
\begin{lstlisting}[language=bash]
python3 codex_heat1d_qutip.py
\end{lstlisting}

\subsection{Trajectory CSV + JSON}
\begin{lstlisting}[language=bash]
python3 codex_heat1d_qutip.py \
  --trajectory-steps 80 \
  --trajectory-csv results_bosonic/codex_heat1d_qutip_trajectory.csv \
  --output-json results_bosonic/codex_heat1d_qutip_metrics.json
\end{lstlisting}

\subsection{CHO-style component plots}
Two components:
\begin{lstlisting}[language=bash]
python3 codex_heat1d_qutip.py \
  --trajectory-steps 80 \
  --component-plot results_bosonic/codex_heat1d_qutip_CHOstyle_components.png \
  --plot-components 0,1
\end{lstlisting}

All four components:
\begin{lstlisting}[language=bash]
python3 codex_heat1d_qutip.py \
  --trajectory-steps 80 \
  --component-plot results_bosonic/codex_heat1d_qutip_CHOstyle_all4.png \
  --plot-components 0,1,2,3
\end{lstlisting}

\subsection{High-Fock stress test}
The current coefficient pipeline can become unstable if \texttt{n-coeff}
is too large at high \texttt{n-fock}. A stabilized run uses smaller
\texttt{n-coeff} than \texttt{n-fock}:
\begin{lstlisting}[language=bash]
python3 codex_heat1d_qutip.py --n-fock 256 --n-coeff 120
python3 codex_heat1d_qutip.py --n-fock 300 --n-coeff 140
\end{lstlisting}

\section{Limitations and Scope}
\begin{enumerate}[itemsep=3pt]
  \item This is a dense-matrix QuTiP simulation for correctness checks,
  not a fault-tolerant gate-count resource estimate.
  \item The default \texttt{position-scale=0.8} is empirical.
  \item Best-scale metrics can mask amplitude-calibration issues;
  eta-calibrated metrics should be interpreted separately.
  \item Large-$n$ coefficient generation may require tailored stabilization
  (e.g., lower \texttt{n-coeff}, modified quadrature strategy).
\end{enumerate}

\section{Conclusion}
The script provides a transparent, reproducible CV-DV LCHS heat-equation
reference with explicit operator-level and state-level diagnostics. In the
current implementation, baseline accuracy is good at moderate truncation
(\texttt{n-fock=64}) and trajectory behavior is consistent with the reported
metrics. The document and commands here are sufficient for exact reproduction
of all generated artifacts in this repository.

\end{document}
